% Generated draft: requires mathematical and source review.
\documentclass[11pt]{article}
\usepackage[margin=25mm]{geometry}
\usepackage{fontspec}
\setmainfont{Latin Modern Roman}
\usepackage{amsmath,amssymb,mathtools,unicode-math}
\setmathfont{Latin Modern Math}
\usepackage{enumitem,needspace,xurl,hyperref}
\setlength{\parindent}{0pt}
\setlength{\parskip}{5pt}
\setlength{\emergencystretch}{3em}
\newcommand{\N}{\mathbb{N}}
\newcommand{\Z}{\mathbb{Z}}
\newcommand{\Q}{\mathbb{Q}}
\newcommand{\R}{\mathbb{R}}
\newcommand{\C}{\mathbb{C}}
\newcommand{\F}{\mathbb{F}}
\newcommand{\A}{\mathbb{A}}
\newcommand{\PP}{\mathbb{P}}
\newcommand{\E}{\mathbb{E}}
\newcommand{\eps}{\varepsilon}
\newcommand{\dd}{\,\mathrm{d}}
\begin{document}
\section*{UnsolvedMath ID: 10900048}
\textbf{Source problem code:} catalogue entry
\subsection{Definitions and mathematical statement}
In the standard Heegaard setting let \(S\) be a closed oriented genus-\(g\geq2\) surface, let \(x\in\mathcal T(S)\), and use the Teichm\"uller metric \(d_T\). Write \(\Gamma=\operatorname{Mod}(S)\). A mapping-class counting interpretation of the source's orbit-ball sample is the finite set
\[A_r(x)=\{\varphi\in\Gamma:d_T(x,\varphi x)\leq r\}.\]
Let \(P\subset\Gamma\) consist of pseudo-Anosov classes whose invariant measured foliations have the stated generic trivalent singularities (three-pronged singularities in the unpunctured surface setting). Fix the handlebodies and boundary markings defining a genus-\(g\) Heegaard gluing, including a fixed orientation-reversing identification, and write \(M_\varphi\) for the resulting closed manifold. Put
\[H=\{\varphi\in\Gamma:M_\varphi\text{ admits a hyperbolic metric}\}.\]
Maher's compound assertion, with this sample convention, is
\[\lim_{r\to\infty}\frac{|A_r(x)\cap P|}{|A_r(x)|}=1,\qquad
\lim_{r\to\infty}\frac{|A_r(x)\cap H|}{|A_r(x)|}=1.\]
Both conclusions are required; the first includes trivalent singularities rather than only pseudo-Anosov genericity. The second concerns gluing the fixed handlebodies, not an unspecified three-manifold associated to a mapping class.

The source literally counts orbit points \(\varphi x\), whereas pseudo-Anosov type and gluing are properties of \(\varphi\). Distinct mapping classes giving the same point differ by its stabilizer, and these properties are not specified as constant on that coset. To count distinct orbit points instead of mapping classes, one must choose representatives or justify descent of both properties. Thus the displayed group-element interpretation is explicit, not a silent correction. The surface type, genus, choice/quantifier of \(x\), stabilizer convention, handlebody markings and meaning of generic singularities must be fixed from context. Punctured or low-genus variants require different singularity/Heegaard conventions.

The sample here is a Teichm\"uller displacement ball with uniform finite counting. A random-walk result or a word-ball result uses a different distribution and does not alone establish either displayed limit.
\subsection{Short English statement}
Uniform counts in large Teichm\"uller orbit balls should be dominated both by pseudo-Anosov maps with trivalent foliations and by hyperbolic Heegaard gluings. Orbit-point stabilizers and the fixed gluing data must be specified.
\subsection{Sources}
\begin{itemize}
\item Local UnsolvedMath registry, permanent ID \(10900048\), exact source code \texttt{catalogue entry}.
\item Maher, attribution in Delp, Hoffoss and Manning, \emph{Problems In Groups, Geometry, and Three-Manifolds} (2015), subsection 6.4, PDF page 5; both assertions share this source label: \url{https://arxiv.org/pdf/1512.04620#page=5}; abstract \url{https://arxiv.org/abs/1512.04620}. Registry access 2026-07-29; source-tex extraction recorded as checked against rendered PDF text.
\item The summary mentions Maher's random-walk work, the virtually Haken/hyperbolization programme and a parallel registry item 0045, without exact theorem citations or the content of that item.
\end{itemize}
\end{document}
